\honest{The Lens That Sees Its Flaws}{Eliezer Yudkowsky}{2007}

Light leaves the Sun, bounces off your shoelaces, enters your eye, triggers nerve impulses, and is rebuilt in the brain as a model of an untied shoelace; so you believe your shoelaces are untied. None of this is magic, and you can understand it, and so think about which ways of thinking make beliefs mirror reality and which do not.

A mouse can see but cannot understand seeing. Its world holds cats, holes and cheese but not mouse brains, so it cannot correct for an optical illusion: its camera takes no pictures of its own lens. We can look at a strange image and realize that part of what we see is the lens. You need not always believe your eyes, but you must know that you have eyes, and keep separate places in your mind for the map and the territory. I ask you to marvel at this, since it is rare among animals.

Science is this same understanding made systematic: reasoning about a more reliable process for making the mind mirror the world. Thinking about ``performing replicable experiments to falsify theories,'' we can see why it works. It is not a separate realm for laboratories but an understandable process that ties brains to reality. The claim I actually need, that a mind can be understood and repaired the way a machine can, gets one sentence. I defend it by saying that the mind may be ``a little more complicated'' than a steam engine. \nb{That is a joke, and it is the whole defense.}

Now my one example. Years ago, in a chatroom, I asked strangers whether they expected a nuclear war in the next twenty years, and asked the ones who said no to explain themselves. \nb{The ones who said yes were not asked to explain anything.} One man said he expected no war for a hundred years, because ``All of the players involved in decisions regarding nuclear war are not interested right now.'' I asked why he stretched it to a hundred years, and he said ``Pure hope.'' From those two words I work out that the thought of nuclear war made him unhappy, and that his brain therefore rejected it.

To show that hoping for peace does not make peace likelier, I bring in parallel universes and a paper by Max Tegmark. \nb{Ordinary probability would have done the job.} Across a billion such worlds, his way of thinking would not make optimists more common in the worlds without war. The principle is that wanting something to be true is not evidence that it is true. Asking which beliefs make you happy looks inward and tells you about yourself; happiness should follow from your picture of the world, not shape it. \nb{People call this wishful thinking, and have for a long time.}

I conclude that a mind which knows its own flaws can apply a reflective correction: if you see that hope is shifting your beliefs too far, you can correct for it. A lens that understands its own systematic errors is not perfect but ``far more powerful.'' I do not say how you would know the correction had worked. \nb{Five months earlier, in ``Knowing About Biases Can Hurt People,'' Yudkowsky had written that knowing about biases can make people reason worse. This post does not mention that.}

A footnote answers a clever objection: hope makes me work harder, which helps the economy and so makes war less likely. We must bring in Bayes's theorem and measure the effect, and your optimism cannot lower the chance of nuclear war by 20\%, or by however much it moved your belief. I give no other number. \nb{The footnote's point, that a small effect should move your beliefs only a little, is the most useful thing in the essay.}
